% TOP_PROBLEM: {"id": 3096, "title": "The Hodge Conjecture", "category_id": 6, "category": {"id": 6, "name": "geometry", "display_name": "Geometry", "description": "Euclidean and non-Euclidean geometry, geometric structures.", "slug": "geometry", "order_index": 6, "created_at": "2026-07-31T15:26:25.671Z"}, "status": "open", "problem_number": "OPG-1803", "set_id": 12, "statusQualification": "The standard rational Hodge conjecture is used, explicitly restoring smoothness of the ambient projective variety."}
% FIRST_POSTED: 2026-10-01
% ATTEMPT_STATUS: unresolved
% COMPLETION: 5
% MODEL: GPT 6 Astra Ultra
% UnsolvedMath ID 3096; source catalogue code OPG-1803
% !TeX program = lualatex
% Research attempt, 2026-10-01. Canonical statement preserved verbatim.
\documentclass[11pt,a4paper]{article}
\usepackage[margin=25mm]{geometry}
\usepackage{fontspec}
\setmainfont{lmroman10-regular.otf}[BoldFont=lmroman10-bold.otf,
 ItalicFont=lmroman10-italic.otf,BoldItalicFont=lmroman10-bolditalic.otf]
\usepackage{amsmath,amssymb,mathtools}
\usepackage{unicode-math}
\setmathfont{latinmodern-math.otf}
\AtBeginDocument{\let\setminus\smallsetminus}
\usepackage{xurl,hyperref}
\hypersetup{colorlinks=true,urlcolor=blue}
\setcounter{secnumdepth}{0}
\setlength{\parindent}{0pt}
\setlength{\parskip}{5pt}
\setlength{\emergencystretch}{3em}
\begin{document}
\section*{The Hodge Conjecture}\label{3096}
{\small UnsolvedMath ID 3096 · OPG-1803 · Geometry · OpenGarden}
\subsection{Definitions and mathematical statement}
Let $X$ be a smooth projective algebraic variety over $\mathbb C$ of
complex dimension $d$. Thus $X$ is a nonsingular closed algebraic subvariety
of some complex projective space. For $0\le p\le d$, its singular
cohomology with complex coefficients has Hodge decomposition
\[
 H^{2p}(X,\mathbb C)=\bigoplus_{a+b=2p}H^{a,b}(X).
\]
The summand $H^{a,b}(X)$ is represented by harmonic differential forms
of complex bidegree $(a,b)$, using a Kähler metric.
Via extension of coefficients, view $H^{2p}(X,\mathbb Q)$ as a rational
subspace of $H^{2p}(X,\mathbb C)$. A rational Hodge class of codimension
$p$ is an element
$\alpha\in H^{2p}(X,\mathbb Q)\cap H^{p,p}(X)$.

An irreducible closed algebraic subvariety $Z\subseteq X$ of complex
codimension $p$ has a fundamental cycle, whose Poincaré dual gives a
cohomology class $[Z]\in H^{2p}(X,\mathbb Q)$. The subvariety $Z$ itself
need not be smooth.
The rational Hodge conjecture asserts that every rational Hodge class
can be expressed as
\[
 \alpha=\sum_{j=1}^m q_j[Z_j],\qquad q_j\in\mathbb Q,
\]
for a finite list of codimension-$p$ algebraic subvarieties of $X$.
All dimensions, all $p$ in the indicated range, and all such classes
are quantified over.

Smoothness of $X$ is an essential scope condition omitted from the short
catalogue statement and supplied here in the classical formulation.
Rational coefficients are essential too: the assertion neither demands
a single effective subvariety nor requires integral coefficients in the
linear combination.
\subsection{Short English statement}
On a smooth complex projective variety, is every rational cohomology class of type $(p,p)$ a rational combination of algebraic subvariety classes?
\subsection{Research attempt}
\paragraph{Scope and literature check.}
Deligne's official description [S3], checked 1 October 2026, is the reference
for rational coefficients and algebraic cycle classes. This attempt tests a
product-based strategy. It proves the assertion for products of projective
spaces, then shows concretely why combining only classes pulled back from
factors fails even on a product of elliptic curves. The standard cohomology
ring of projective space and the rational K\"unneth formula are used as
established inputs [A1], not as conclusions about a general projective variety.

\paragraph{A complete restricted calculation.}
Let $X=\prod_{i=1}^r\mathbb P^{d_i}_{\mathbb C}$ and let $h_i$ be the pullback
of the hyperplane class from factor $i$. The projective-space cell
decomposition and the K\"unneth formula give
\[
 H^*(X,\mathbb Q)=\mathbb Q[h_1,\ldots,h_r]/
                   (h_1^{d_1+1},\ldots,h_r^{d_r+1}),\qquad \deg h_i=2.
\]
The monomials $h_1^{a_1}\cdots h_r^{a_r}$ with $0\le a_i\le d_i$ and
$\sum_i a_i=p$ form a rational basis of $H^{2p}(X,\mathbb Q)$.
Each is the class of the algebraic subvariety
\[
                  \prod_i\mathbb P^{d_i-a_i}\subseteq X,
\]
obtained by choosing linear subspaces in the factors. Indeed intersecting
$a_i$ independent hyperplanes in factor $i$ gives its displayed linear
subspace with multiplicity one, and products of these intersections give
the cup product class. Thus every rational degree-$2p$ class, and in
particular every rational Hodge class, is a rational sum of algebraic cycles.
This proves the catalogue assertion for this entire family in every
codimension. The coefficients may be negative; requiring effectivity would
change the assertion and is not part of the proof.

\paragraph{A sharp obstruction to the naive product extension.}
Let $E$ be a complex elliptic curve, $X=E\times E$, and choose $o\in E$.
Put $F_1=E\times\{o\}$, $F_2=\{o\}\times E$, and let $\Delta$ be the diagonal.
The basic intersection numbers on this smooth surface are
\[
 F_1^2=F_2^2=\Delta^2=0,\qquad F_1F_2=\Delta F_1=\Delta F_2=1.
\]
For the first two self-intersections, translate the fixed point to a distinct
point to obtain a disjoint homologous fibre. For $\Delta^2$, a graph
$\{(x,x+t):x\in E\}$ with $t\ne0$ is disjoint from the diagonal and homologous
to it by a continuous family of translations. The remaining intersections
are transverse at one point, hence have multiplicity one.
If $[\Delta]=a[F_1]+b[F_2]$, intersecting with the fibres forces $a=b=1$.
Squaring then gives $\Delta^2=2$, contradicting zero. Therefore even this
algebraic Hodge class is not in the span of the degree-two classes pulled
back from the two factors.

\paragraph{What the countercalculation diagnoses.}
For projective spaces, there is no odd cohomology and all even cohomology is
generated by hyperplanes. For $E\times E$, K\"unneth includes the additional
summand $H^1(E,\mathbb Q)\otimes H^1(E,\mathbb Q)$ inside degree two.
The diagonal has a nonzero component there, as the intersection calculation
proves. Thus an induction that only multiplies factorwise even-dimensional
cycle bases omits actual algebraic Hodge classes. The example does not
refute the Hodge conjecture: $\Delta$ is itself a perfectly good algebraic
cycle. It refutes a proposed spanning argument that is too narrow.

\paragraph{Verification and remaining gap.}
The projective-space proof handles $p=0$ and top degree, giving the whole
variety and a point respectively. The elliptic calculation uses rational
cohomology and intersection pairings, not an invalid conclusion about
integral Hodge classes. A continuation would need to represent mixed Hodge
classes by correspondences or other subvarieties and prove that they span
all such classes. No construction achieving that for arbitrary $X$ is given.

\paragraph{Final status.}
\textbf{UNRESOLVED.} A standard positive family has been proved explicitly
and a tempting factorwise extension has been disproved. Neither calculation
is claimed as new, and neither settles the general Hodge conjecture.
Completion estimate: 5\%.

\subsection{Sources}
\begin{itemize}
\item[S1] ULAM.AI and the UnsolvedMath Contributors, supplied problems.json, record 3096 (OPG-1803), OpenGarden collection. Catalogue identity and source transcription; CC BY 4.0.
\url{https://huggingface.co/datasets/ulamai/UnsolvedMath}.
\item[S2] Open Problem Garden, The Hodge Conjecture. Original problem and discussion; the mathematical formulation below is expanded from the supplied source record.
\url{http://www.openproblemgarden.org/op/the_hodge_conjecture}.
\item[S3] P. Deligne, The Hodge Conjecture, official problem description, Clay Mathematics Institute.
\url{https://www.claymath.org/wp-content/uploads/2022/06/hodge.pdf}.

\item[A1] A. Hatcher, \emph{Algebraic Topology}, Chapter 3 (cup products, projective-space cohomology and the K\"unneth formula). Author's book site checked 1 October 2026. \url{https://pi.math.cornell.edu/~hatcher/AT/ATpage.html}.
\end{itemize}
\end{document}
